\documentclass[12pt,article,a4paper,brazil]{abntex2}

% --- 
% CONFIGURAÇÕES DE PACOTES
% --- 

\usepackage[T1]{fontenc}              % seleção de códigos de fonte.
\usepackage[utf8]{inputenc}              % determina a codificação utiizada (conversão automática dos acentos)
\usepackage[brazil]{babel}                  % idiomas
\usepackage{hyperref}                    % controla a formação do índice
\usepackage{parskip}
\usepackage[alf]{abntex2cite}                 % Citações padrão ABNT
\usepackage[brazilian,hyperpageref]{backref}  % indique quantas vezes e em quais páginas a citação ocorreu
\usepackage{lipsum}
\usepackage{wallpaper}
\usepackage{graphicx}
\usepackage{listings}
\usepackage{graphicx}
\usepackage[colorinlistoftodos]{todonotes}
\usepackage{indentfirst}
\usepackage[colorinlistoftodos]{todonotes}



% ---
% Configurações do pacote backref
% Usado sem a opção hyperpageref de backref
\renewcommand{\backrefpagesname}{Citado na(s) página(s):~}
% Texto padrão antes do número das páginas
\renewcommand{\backref}{}
% Define os textos da citação
\renewcommand*{\backrefalt}[4]{
	\ifcase #1 %
		Nenhuma citação no texto.%
	\or
		Citado na página #2.%
	\else
		Citado #1 vezes nas páginas #2.%
	\fi}%
% ---


\titulo{UFJF -- IAD Template de artigo v.0.0.2}

\autor{Primeiro autor
  \url{primeiro@email.net}
  \thanks{Titulação Acadêmica} 
  \and 
  Segundo Autor
  \url{segundo@email}
  \thanks{Titulação} 
}

\data{\today}
\thispagestyle{empty}
\begin{document}

% Adiciona imagens da instituição
% TODO Ajustar posições
\ThisULCornerWallPaper{0.2}{ufjf} %size, name
\ThisURCornerWallPaper{0.2}{logoIAD} %size, name

\maketitle
\pagenumbering{arabic}

\begin{abstract}
Coloque seu resume aqui; não ultrapasse 800 caracteres.
\end{abstract}

\section{INTRODUÇÃO}

\lipsum{1}

\section{Alguns Exemplos \LaTeX e \abnTeX}
\label{sec:examples}

\subsection{COMO DEIXAR COMENT\'ARIOS}

\lipsum \todo{Here's a comment in the margin!} \lipsum

\todo[inline, color=green!40]{This is an inline comment.}

\subsection{How to Include Figures}

\lipsum

\begin{figure}
\centering
\includegraphics[width=0.3\textwidth]{frog.jpg}
\caption{\label{fig:frog}This frog was uploaded to writeLaTeX via the project menu.}
\end{figure}

\subsection{COMO FAZER LISTAS}

Use the table and tabular commands for basic tables --- see Table~\ref{tab:widgets}, for example.

\begin{table}
\centering
\begin{tabular}{l|r}
Item & Quantity \\\hline
Widgets & 42 \\
Gadgets & 13
\end{tabular}
\caption{\label{tab:widgets}An example table.}
\end{table}

\subsection{COMO FAZER SEÇ\~OES E SUBSEÇ\~OES}

Use section and subsection commands to organize your document. \LaTeX{} handles all the formatting and numbering automatically. Use ref and label commands for cross-references.

\subsection{COMO FAZER LISTAS}

You can make lists with automatic numbering \dots

\begin{enumerate}
\item Like this,
\item and like this.
\end{enumerate}
\dots or bullet points \dots
\begin{itemize}
\item Like this,
\item and like this.
\end{itemize}
\dots or with words and descriptions \dots
\begin{description}
\item[Word] Definition
\item[Concept] Explanation
\item[Idea] Text
\end{description}

We hope you find write\LaTeX\ useful, and please let us know if you have any feedback using the help menu above.

\end{document}